% ============================================================================
% Bild 4: Bemasste "technische Zeichnung" der unendlichen Pol-Kette
% ============================================================================
% MATHEMATISCHER HINTERGRUND / DESIGN-ENTSCHEIDUNGEN:
%
% * Dargestellt werden die Pole s_nu = -nu^2 (nu = 1, 2, 3, ...) der
%   Bildfunktion u_0(x,s) = sinh((pi-x)sqrt(s))/sinh(pi sqrt(s)) des
%   Stabes der Laenge l = pi, sowie die Bogenfolge C_n aus dem Paper
%   mit den im URSPRUNG zentrierten Radien R_n = (n + 1/2)^2, die
%   GENAU ZWISCHEN den Polen hindurchlaufen:
%       R_1 = 2.25   (zwischen -1 und -4)
%       R_2 = 6.25   (zwischen -4 und -9)
%
% * MASSSTAB: 0.4 cm pro s-Einheit (linear bis zur Bruchstelle).
%   Damit liegen die Pole bei x = -0.4, -1.6, -3.6 cm und die
%   Bogenradien betragen 0.9 cm bzw. 2.5 cm.
%
% * BOGEN-GEOMETRIE: Die Boegen schliessen an die Bromwich-Gerade
%   Re(s) = gamma an (gamma = 1 s-Einheit = 0.4 cm), d.h. die Gerade
%   ist SEHNE der Kreise (konsistent zur D-Kontur-Abbildung).
%   Startwinkel aus cos(theta) = 0.4/R:
%       C_1: R = 0.9  -> theta = 63.61 deg, Endpunkt-y = 0.9*sin = 0.806
%       C_2: R = 2.5  -> theta = 80.79 deg, Endpunkt-y = 2.5*sin = 2.468
%
% * BEMASSUNG: Massketten unterhalb der Achse zeigen die Abstaende
%   1, 3, 5 (in s-Einheiten) zwischen 0,-1,-4,-9 -- die ungeraden
%   Zahlen machen das quadratische Wachstum |s_nu| = nu^2 sichtbar.
%   Technischer Stil: duenne graue Hilfslinien + Masslinien mit
%   Endstrichen (Bar-Tips statt Pfeilen, weil das kuerzeste Mass mit
%   0.4 cm zu schmal fuer beidseitige Pfeilspitzen ist).
%
% * BRUCHSTELLE: Doppelter Zickzack-Bruch (techn. Zeichnungsnorm)
%   nach -9; dahinter ein letzter Pol "-n^2" und ein Achsenstummel
%   mit Pfeil nach links (Fortsetzung der Kette bis -unendlich).
%
% * Bei s = 0 KEIN Pol (hebbare Singularitaet, Zaehler verschwindet
%   mit): offener Kreis statt Kreuz, graue Annotation "hebbar".
% ============================================================================
\scalebox{0.85}{
\begin{tikzpicture}

% ==============================
% 1. FARBEN DEFINIEREN
% ==============================
\definecolor{plotblue}{RGB}{0, 0, 200}
\definecolor{plotred}{RGB}{220, 30, 30}
\definecolor{plotgray}{RGB}{140, 140, 140}
\definecolor{plotbg}{RGB}{244, 244, 252}

% ==============================
% 2. BOEGEN C_n (zuunterst, damit alles andere darueber liegt)
% ==============================
% C_1: R = 0.9 cm, Sehne bei x = 0.4  ->  arc von 63.61 bis 296.39 Grad
\draw[thick, plotgray, dashed] (0.4, 0.806) arc (63.61:296.39:0.9);
% C_2: R = 2.5 cm, Sehne bei x = 0.4  ->  arc von 80.79 bis 279.21 Grad
\draw[thick, plotgray, dashed] (0.4, 2.468) arc (80.79:279.21:2.5);

% Bogen-Labels (jeweils knapp ausserhalb des Bogens, oben links)
\node[plotgray, font=\normalsize] at (-0.55, 1.15) {$C_1$};
\node[plotgray, font=\normalsize] at (-1.75, 2.15) {$C_2$};

% ==============================
% 3. BROMWICH-GERADE (Sehne der Boegen bei Re(s) = gamma)
% ==============================
\draw[thick, plotblue, decoration={markings, mark=at position 0.6 with {\arrow{Stealth[length=3mm, width=2mm]}}}, postaction={decorate}]
    (0.4, -2.468) -- (0.4, 2.468);
\draw[thick] (0.4, 0.1) -- (0.4, -0.1) node[below right, inner sep=2pt] {$\boldsymbol{\gamma}$};

% ==============================
% 4. ACHSEN MIT BRUCHSTELLE
% ==============================
% Hauptachse (rechts der Bruchstelle)
\draw[thick, -{Stealth[length=3.5mm, width=2.5mm]}] (-3.9, 0) -- (2.3, 0) node[right] {\textbf{Re}$(s)$};
% Achsenstummel links der Bruchstelle, Pfeil nach links (bis -unendlich)
\draw[thick, -{Stealth[length=3.5mm, width=2.5mm]}] (-4.34, 0) -- (-6.05, 0);
% Doppelter Zickzack-Bruch (techn. Norm fuer "hier fehlt ein Stueck")
\draw[thick] (-3.92, -0.3) -- (-4.04, -0.06) -- (-3.96, 0.06) -- (-4.08, 0.3);
\draw[thick] (-4.14, -0.3) -- (-4.26, -0.06) -- (-4.18, 0.06) -- (-4.30, 0.3);
% Kurze Im-Achse
\draw[thick, -{Stealth[length=3.5mm, width=2.5mm]}] (0, -2.9) -- (0, 3.05) node[above] {\textbf{Im}$(s)$};

% ==============================
% 5. POLE (Kreuze) UND HEBBARE STELLE
% ==============================
\newcommand{\pole}[3][plotred]{
    \draw[thick, #1, line cap=round] (#2-0.12,#3-0.12) -- (#2+0.12,#3+0.12);
    \draw[thick, #1, line cap=round] (#2-0.12,#3+0.12) -- (#2+0.12,#3-0.12);
}

% Pole bei s = -1, -4, -9 (Massstab 0.4 cm/Einheit) und -n^2
\pole{-0.4}{0}
\pole{-1.6}{0}
\pole{-3.6}{0}
\pole{-5.0}{0}

% Pol-Beschriftungen OBERHALB der Achse (unten laeuft die Masskette)
\node[anchor=south, font=\small] at (-0.4, 0.22) {$-1$};
\node[anchor=south, font=\small] at (-1.6, 0.22) {$-4$};
\node[anchor=south, font=\small] at (-3.6, 0.22) {$-9$};
\node[anchor=south, font=\small] at (-5.0, 0.22) {$-n^2$};

% Hebbare Singularitaet bei s = 0: roter Kreis mit rotem Kreuz
\draw[thick, plotred] (0,0) circle (3.0pt);
\draw[thick, plotred, line cap=round] (-0.075,-0.075) -- (0.075,0.075);
\draw[thick, plotred, line cap=round] (-0.075,0.075) -- (0.075,-0.075);



% ==============================
% 7. LEGENDE
% ==============================
\node[draw, thick, rounded corners=3pt, fill=white, inner xsep=5pt, inner ysep=4pt, anchor=north west] at (1.3, 3.05) {
    \renewcommand{\arraystretch}{1.2}
    \begin{tabular}{@{}c l@{}}

        \tikz[baseline=-0.6ex]{
            \draw[thick, plotred, line cap=round] (-0.1,-0.1) -- (0.1,0.1) (-0.1,0.1) -- (0.1,-0.1);
        } & Polstelle \\

        \tikz[baseline=-0.6ex]{
            \draw[thick, plotred] (0,0) circle (3.0pt);
            \draw[thick, plotred, line cap=round] (-0.075,-0.075) -- (0.075,0.075);
            \draw[thick, plotred, line cap=round] (-0.075,0.075) -- (0.075,-0.075);
        } & Hebbare Singularit\"at \\

        \tikz[baseline=-0.6ex]{
            \draw[thick, plotblue, ->-=0.6] (-0.3, 0) -- (0.3, 0);
        } & Bromwich-Pfad \\

        \tikz[baseline=-0.6ex]{
            \draw[thick, plotgray, dashed] (-0.3, 0) -- (0.3, 0);
        } & Bogen $C_n$ \\

    \end{tabular}
};

\end{tikzpicture}
}
